\section{量、单位、测量与数量级}

如果读者从未拆过电脑，也没有学过电路，本书仍然应该能够独立阅读。“从零”
并不表示用比喻替代事实，而是遵守以下顺序：

\[
\text{能观察什么}
\rightarrow
\text{怎样测量}
\rightarrow
\text{怎样表示}
\rightarrow
\text{规则怎样改变状态}
\rightarrow
\text{怎样验证结论}
\]

看到一个新名词时，先不要问“考试定义是什么”，而要问：它是一个物体、一种
可测量的量、一组编码规则、一次状态变化，还是人为划定的抽象边界？例如：

\begin{osgridlongtable}{L{0.18\linewidth}L{0.28\linewidth}L{0.44\linewidth}}
词语 & 类别 & 可以直接追问的问题 \\ \hline
\endfirsthead
词语 & 类别 & 可以直接追问的问题 \\ \hline
\endhead
电阻 & 物理元件 & 两端怎样连接，阻值、功率和误差是多少？\\ \hline
电压 & 可测量的量 & 哪两个点之间，使用什么参考，单位是什么？\\ \hline
bit & 可区分状态的编码 & 哪个物理状态代表 0，哪个代表 1？\\ \hline
地址 & 编号规则 & 编号属于哪个地址空间，由谁译码？\\ \hline
指令 & 状态转换规则 & 输入状态、输出状态和失败条件是什么？\\ \hline
进程 & 软件抽象边界 & 它拥有什么，谁创建和回收它？\\ \hline
\end{osgridlongtable}

\begin{fromzerobox}
“地是零伏”“1 是高电平”“地址就是内存位置”都省略了限定条件。更准确的
说法分别是：电压相对于选定参考点测量；逻辑输入在规定阈值范围内被解释为
高电平；某个地址空间中的编号经过相应译码后选择一个目标。
\end{fromzerobox}

\topic{物体、现象、量、单位和数值不是一回事}
一块电路板是物体；电流流动、芯片发热和时钟振荡是现象；电压、电流、温度
和频率是可以测量的量；伏特、安培、摄氏度和赫兹是单位；\(\SI{3.3}{V}\)、
\(\SI{20}{mA}\) 和 \(\SI{100}{MHz}\) 才是带单位的数值。

一个完整测量至少包含：

\[
Q = q \times [Q]
\]

其中 \(Q\) 是物理量，\(q\) 是数值，\([Q]\) 是单位。只写“电压是 3.3”
会丢失单位；只写“这个点是 3.3 V”仍会丢失参考点和工作条件。

\topic{七个常见量先建立直觉}
\begin{osgridlongtable}{L{0.16\linewidth}L{0.12\linewidth}L{0.19\linewidth}L{0.43\linewidth}}
量 & 符号 & 常用单位 & 在计算机硬件中的问题 \\ \hline
\endfirsthead
量 & 符号 & 常用单位 & 在计算机硬件中的问题 \\ \hline
\endhead
时间 & \(t\) & s、ms、\(\mu\)s、ns & 一个周期多长，超时等待多久？\\ \hline
频率 & \(f\) & Hz、MHz、GHz & 每秒重复多少周期？\\ \hline
电压 & \(V\) & V、mV & 两点电势差是否落在器件允许范围？\\ \hline
电流 & \(I\) & A、mA & 某支路每秒通过多少电荷？\\ \hline
电阻 & \(R\) & \(\Omega\)、k\(\Omega\) & 给定电压会形成多大电流？\\ \hline
功率 & \(P\) & W、mW & 单位时间转换多少能量，会不会过热？\\ \hline
数据量 & \(N\) & bit、B、KiB、MiB & 能表示多少状态、保存多少字节？\\ \hline
\end{osgridlongtable}

周期与频率互为倒数：

\[
T=\frac{1}{f}
\]

若时钟频率为 \(\SI{100}{MHz}=100\times10^6\ \mathrm{s}^{-1}\)，周期为

\[
T=\frac{1}{100\times10^6}\ \mathrm{s}
=10^{-8}\ \mathrm{s}
=\SI{10}{ns}.
\]

这不表示 CPU 在每个周期必然“完成一条指令”。频率只描述节拍；指令吞吐还
取决于流水线、依赖、缓存、分支和微架构。

\topic{前缀每差三级就是一千倍}
\begin{center}
\begin{osgridtabular}{llll}
前缀 & 符号 & 倍率 & 示例 \\ \hline
giga & G & \(10^9\) & \(\SI{3}{GHz}\) \\ \hline
mega & M & \(10^6\) & \(\SI{100}{MHz}\) \\ \hline
kilo & k & \(10^3\) & \(\SI{10}{k\ohm}\) \\ \hline
milli & m & \(10^{-3}\) & \(\SI{20}{mA}\) \\ \hline
micro & \(\mu\) & \(10^{-6}\) & \(\SI{5}{\micro s}\) \\ \hline
nano & n & \(10^{-9}\) & \(\SI{10}{ns}\) \\ \hline
pico & p & \(10^{-12}\) & \(\SI{22}{pF}\) \\ \hline
\end{osgridtabular}
\end{center}

大小写有意义：\(\mathrm{m}\) 是 milli，\(\mathrm{M}\) 是 mega；二者相差
\(10^9\) 倍。工程记录中应写 \(\SI{10}{k\ohm}\)，不写含糊的“10K”。

\section{怎样记录一次可以复查的测量}
真正的测量包含被测对象、被测量、参考点、仪器、连接方法、量程、工作条件和
读数。缺少其中任何一项，数字都可能无法复现。例如“复位脚是 1”只是一种逻辑
解释；“上电后 \(\SI{18}{ms}\)，示波器探头相对板上数字地测得 RESET\_N 为
\(\SI{3.26}{V}\)”才接近一条可复查记录。

\begin{pdfdiagram}[node distance=7mm and 11mm]
  \node[os hardware,minimum width=25mm] (object) {被测对象\\电路节点或事件};
  \node[os block,right=of object,minimum width=24mm] (quantity) {选择量与参考\\电压、时间、频率};
  \node[os hardware,right=of quantity,minimum width=24mm] (instrument) {仪器与连接\\量程、探头、带宽};
  \node[os state,right=of instrument,minimum width=22mm] (reading) {原始读数\\数值、单位、波形};
  \node[os block,below=9mm of instrument,minimum width=32mm] (judgement) {结合误差和规格判断\\支持或推翻一个结论};
  \draw[os arrow] (object) -- (quantity);
  \draw[os arrow] (quantity) -- (instrument);
  \draw[os arrow] (instrument) -- (reading);
  \draw[os arrow] (reading) |- (judgement);
  \draw[os arrow] (quantity) |- (judgement);
\end{pdfdiagram}
\diagramnote{仪器不会直接给出“系统正确”。它只在规定连接和误差范围内提供
观测，再由规格和推理把观测转成结论。}

\topic{基本量、导出量和量纲检查}
有些量作为计量体系的基础；另一些量由关系式导出。本书常用的关系包括：

\[
\text{频率}\ [f]=\mathrm{s}^{-1},\qquad
\text{电荷}\ [Q]=\mathrm{A\,s},\qquad
\text{电压}\ [V]=\mathrm{J/C},\qquad
\text{功率}\ [P]=\mathrm{J/s}.
\]

\term{量纲检查}{dimensional analysis}不能证明公式一定正确，却能迅速发现
许多错误。等号两边必须描述同一种量。例如 \(T=1/f\) 的右侧单位为
\(\mathrm{Hz}^{-1}=\mathrm{s}\)，与周期一致；若有人写 \(T=f/2\)，右侧仍是
\(\mathrm{s}^{-1}\)，仅看单位就知道不可能是时间。

\begin{workedexample}
一条串行链路每个帧发送 10 bit，其中 8 bit 是数据，波特率为
\(\SI{115200}{bit/s}\)。每秒最多完成的帧数为

\[
\frac{\SI{115200}{bit/s}}{\SI{10}{bit/frame}}
=\SI{11520}{frame/s}.
\]

每帧承载 1 byte，因此有效数据率为
\(\SI{11520}{B/s}\)，不是 \(\SI{115200}{B/s}\)。分母中的
\(\mathrm{bit/frame}\) 与分子中的 bit 消去后留下 frame/s，这也同时检查了
推导的单位。
\end{workedexample}

\section{科学计数法与数量级}
计算机硬件同时处理 GHz、ns、mV、pF 和 GiB。科学计数法把非零数写成

\[
a\times10^n,\qquad 1\le |a|<10.
\]

乘法时有效数字相乘、指数相加；除法时指数相减：

\[
(3.3\times10^0)(2.0\times10^{-3})
=6.6\times10^{-3},
\]
\[
\frac{5.0\times10^9}{2.5\times10^8}
=2.0\times10^1.
\]

工程前缀通常让十进制指数是 3 的倍数。把
\(\SI{0.0000047}{F}\) 写成 \(\SI{4.7}{\micro F}\)，不是改变数值，而是选择
更便于阅读的尺度。

\begin{workedexample}
把 \(\SI{25}{ns}\) 换算为秒：

\[
\SI{25}{ns}=25\times10^{-9}\ \mathrm{s}
=2.5\times10^{-8}\ \mathrm{s}.
\]

一秒中包含的这类时间片数量为

\[
\frac{\SI{1}{s}}{\SI{25}{ns}}
=\frac{1}{25\times10^{-9}}
=4.0\times10^7.
\]

所以周期 \(\SI{25}{ns}\) 对应频率 \(\SI{40}{MHz}\)。换算时保留单位可以
避免把 \(10^{-9}\) 错写成 \(10^9\)。
\end{workedexample}

\section{仪器怎样接入电路}
仪器本身也是电路的一部分。接错位置不仅会得到错误读数，还可能短路电源或
损坏被测板。

\begin{osgridlongtable}{L{0.16\linewidth}L{0.22\linewidth}L{0.25\linewidth}L{0.27\linewidth}}
仪器/模式 & 连接方式 & 主要观察 & 常见错误 \\ \hline
\endfirsthead
仪器/模式 & 连接方式 & 主要观察 & 常见错误 \\ \hline
\endhead
万用表电压档 & 并联在两个节点之间 & 近似稳定的电势差 & 把表笔串进回路，或忘记参考地 \\ \hline
万用表电流档 & 断开支路后串联接入 & 经过该支路的平均电流 & 直接跨接电源，等价于近似短路 \\ \hline
万用表电阻档 & 断电并释放储能后跨接元件 & 低频等效电阻 & 带电测量，或把并联路径误当元件本体 \\ \hline
示波器 & 探头尖端接信号，地夹接允许的参考点 & 电压随时间的变化 & 地夹造成短路，探头带宽/衰减设置错误 \\ \hline
逻辑分析仪 & 多路数字输入与公共参考相连 & 阈值判定后的数字时序 & 把模拟边沿质量问题隐藏成整齐的 0/1 \\ \hline
频率计/计时器 & 在允许电平和带宽内接收周期信号 & 周期、频率或事件间隔 & 输入幅度、耦合方式或触发条件不匹配 \\ \hline
\end{osgridlongtable}

\begin{center}
\begin{tikzpicture}[font=\footnotesize\sffamily,>=Latex]
  \draw[BookMuted,thick] (0,2.4) -- (4.8,2.4);
  \draw[BookMuted,thick] (0,0) -- (4.8,0);
  \node[left] at (0,2.4) {\(V_{\mathrm{DD}}\)};
  \node[left] at (0,0) {GND};
  \node[draw=BookAmber,fill=white,minimum width=15mm,minimum height=12mm]
    (load) at (2.4,1.2) {负载};
  \draw (2.4,2.4) -- (load.north);
  \draw (load.south) -- (2.4,0);

  \node[draw=BookBlue,fill=BookCodeBg,minimum width=16mm,minimum height=11mm]
    (voltmeter) at (6.7,1.2) {电压表 \(V\)};
  \draw[BookBlue,thick] (4.4,2.4) -| (voltmeter.north);
  \draw[BookBlue,thick] (4.4,0) -| (voltmeter.south);
  \node[align=center,BookBlue] at (6.7,-0.45) {并联：比较两个节点};

  \draw[BookMuted,thick] (9.0,2.4) -- (10.0,2.4);
  \node[draw=BookTeal,fill=BookCodeBg,minimum width=16mm,minimum height=11mm]
    (ammeter) at (11.0,2.4) {电流表 \(A\)};
  \draw[BookMuted,thick] (12.0,2.4) -- (13.0,2.4);
  \node[draw=BookAmber,fill=white,minimum width=15mm,minimum height=12mm]
    (load2) at (13.0,1.2) {负载};
  \draw (13.0,2.4) -- (load2.north);
  \draw (load2.south) -- (13.0,0) -- (9.0,0);
  \node[align=center,BookTeal] at (11.0,-0.45) {串联：让支路电流经过仪器};
\end{tikzpicture}
\end{center}
\diagramnote{电压表理想输入电阻无限大，电流表理想内阻为零；真实仪器只能
接近这些理想模型，因此仍会改变被测电路。}

\topic{分辨率、准确度、精密度和采样率不是一回事}
\begin{osgridlongtable}{L{0.18\linewidth}L{0.34\linewidth}L{0.38\linewidth}}
概念 & 回答的问题 & 例子 \\ \hline
\endfirsthead
概念 & 回答的问题 & 例子 \\ \hline
\endhead
分辨率 & 最小能显示或区分多小的变化？ & 三位半万用表在某量程显示到 \(\SI{1}{mV}\) \\ \hline
准确度 & 读数离真实值可能有多远？ & \(\pm(0.5\%\text{ 读数}+2\text{ 字})\) \\ \hline
精密度/重复性 & 重复测量聚集得有多紧？ & 多次读数都在 \(\SIrange{3.298}{3.301}{V}\) \\ \hline
采样率 & 每秒取得多少样本？ & \(\SI{100}{MS/s}\) \\ \hline
模拟带宽 & 多快的变化仍能以规定幅度误差通过？ & \(\SI{20}{MHz}\) 示波器不适合判断 GHz 边沿 \\ \hline
\end{osgridlongtable}

显示很多小数位不代表准确。一个稳定显示
\(\SI{3.3000}{V}\) 的仪器若校准偏差为 \(\SI{50}{mV}\)，精密但不准确；
一个准确度足够但采样率很低的万用表可以看到平均电压，却看不到
\(\SI{5}{ns}\) 毛刺。

\topic{仪器负载效应：观察者会改变被观察对象}
假设信号源可等效为理想电压 \(V_s\) 串联源电阻 \(R_s\)，电压表输入电阻为
\(R_m\)。接表后读数是

\[
V_m=V_s\frac{R_m}{R_s+R_m}.
\]

\begin{workedexample}
\(V_s=\SI{3.3}{V}\)、\(R_s=\SI{1}{M\ohm}\)，表的输入电阻
\(R_m=\SI{10}{M\ohm}\)：

\[
V_m=\SI{3.3}{V}\times\frac{10}{1+10}
=\SI{3.0}{V}.
\]

仪表精度即使完美，也只会诚实报告被它拉低后的 \(\SI{3.0}{V}\)。高阻节点
应使用更高输入阻抗的缓冲器或探头，并把探头电容纳入高速分析。
\end{workedexample}

\section{误差传播与有效数字}
若量 \(x\) 的标称值为 \(x_0\)，绝对误差不超过 \(\Delta x\)，可写成

\[
x=x_0\pm\Delta x.
\]

相对误差为

\[
\varepsilon_x=\frac{\Delta x}{|x_0|}.
\]

对于独立量的粗略最坏情况估算，乘除结果的相对误差可近似相加。若
\(R=\SI{10}{k\ohm}\pm1\%\)，\(C=\SI{100}{nF}\pm10\%\)，则
\(\tau=RC\) 的最坏相对偏差近似为 \(11\%\)。因此把时间常数报告成
\(\SI{1.000000}{ms}\) 会制造并不存在的精度。

\begin{workedexample}
一个 \(\SI{1}{V}\) 标称基准误差 \(\pm0.5\%\)，测量仪器在该点的总误差
\(\pm\SI{8}{mV}\)。若只做保守上界相加：

\[
\Delta V_{\mathrm{ref}}=\SI{1}{V}\times0.005=\SI{5}{mV},
\]
\[
\Delta V_{\mathrm{total}}\le\SI{5}{mV}+\SI{8}{mV}
=\SI{13}{mV}.
\]

读数 \(\SI{1.006}{V}\) 与理想值的差小于误差上界，不能据此宣布电路异常。
测量结论必须与不确定度一起表达。
\end{workedexample}

\section{比例、斜率、面积与对数}
硬件公式不是需要孤立背诵的符号：

\begin{itemize}[leftmargin=2.2em]
  \item \(y=kx\) 表示比例关系，输入加倍时输出也加倍；
  \item \(\mathrm{d}y/\mathrm{d}t\) 是变化率，例如电流是电荷变化率；
  \item 曲线下面积表示累积量，例如功率对时间积分得到能量；
  \item \(2^n\) 表示每增加一个二值选择，组合数量翻倍；
  \item \(\log_2 N\) 反问“至少需要多少个二值位置才能区分 \(N\) 个状态”；
  \item \(e^{-t/\tau}\) 描述许多储能系统向稳态逼近的过程。
\end{itemize}

\begin{center}
\begin{tikzpicture}[x=0.85cm,y=0.85cm,font=\footnotesize\sffamily]
  \draw[->,BookMuted] (0,0) -- (6.2,0) node[right]{输入 \(x\)};
  \draw[->,BookMuted] (0,0) -- (0,4.2) node[above]{输出 \(y\)};
  \draw[BookBlue,very thick] (0,0) -- (5.4,3.6);
  \draw[dashed,BookRule] (1.5,0) -- (1.5,1.0) -- (0,1.0);
  \draw[dashed,BookRule] (4.5,0) -- (4.5,3.0) -- (0,3.0);
  \node[BookBlue] at (4.4,3.55) {\(y=kx\)};
  \node[align=center] at (3.0,-0.75) {斜率 \(k=\Delta y/\Delta x\)\\表示每单位输入带来多少输出};

  \begin{scope}[xshift=8cm]
    \draw[->,BookMuted] (0,0) -- (6.2,0) node[right]{时间 \(t\)};
    \draw[->,BookMuted] (0,0) -- (0,4.2) node[above]{响应};
    \draw[BookTeal,very thick,samples=80,domain=0:5.5,smooth]
      plot (\x,{3.6*(1-exp(-\x/1.3))});
    \draw[dashed,BookRule] (0,3.6) -- (5.7,3.6);
    \draw[dashed,BookAmber] (1.3,0) -- (1.3,2.28);
    \node[BookAmber,below] at (1.3,0) {\(\tau\)};
    \node[align=center] at (3.0,-0.75) {指数响应先快后慢\\在 \(t=\tau\) 到达约 \(63.2\%\)};
  \end{scope}
\end{tikzpicture}
\end{center}

\topic{估算先于精算：数量级能拦住荒谬答案}
在使用计算器前先估算数量级。若 \(\SI{3.3}{V}\) 加在
\(\SI{1}{k\ohm}\) 上，电流应在 mA 量级；若算出 \(\SI{3.3}{kA}\)，说明单位
换算错了。若 \(\SI{1}{GHz}\) 时钟的周期算成一秒，也应立即被常识拦下。

一个实用检查顺序是：

\begin{enumerate}[leftmargin=2.2em]
  \item 先写物理关系，不急着代数字；
  \item 把所有前缀展开成 \(10^n\) 或 \(2^n\)；
  \item 让单位在式子里约掉；
  \item 估算数量级和上下界；
  \item 再计算有效数字；
  \item 用反向代入或另一种方法核对。
\end{enumerate}

\begin{failurebox}
公式计算器只会忠实执行输入。把 \(\SI{100}{nF}\) 当成
\(\SI{100}{F}\)、把 byte 当 bit、把峰峰值当幅值、把十六进制字符串当十进制，
都可能得到形式整齐但物理荒谬的结果。单位、参考点和数量级检查不能省略。
\end{failurebox}


\section{进制、固定位宽与内存布局}

位置计数法把一串数字解释成不同权重的和。十进制数 \(507\) 表示

\[
5\times10^2+0\times10^1+7\times10^0.
\]

二进制 \(101101_2\) 表示

\[
1\times2^5+0\times2^4+1\times2^3+1\times2^2+0\times2^1+1\times2^0
=32+8+4+1=45.
\]

十六进制使用 \(0\)--\(9\) 和 \(A\)--\(F\) 表示 0--15。因为
\(16=2^4\)，一个十六进制数字恰好对应四个 bit：

\begin{center}
\begin{osgridtabular}{cccccccc}
二进制 & 0000 & 0001 & 0010 & 0111 & 1000 & 1010 & 1111 \\ \hline
十六进制 & 0 & 1 & 2 & 7 & 8 & A & F \\ \hline
\end{osgridtabular}
\end{center}

因此

\[
1011\,0110\,1111\,0001_2 = \texttt{0xB6F1}.
\]

十六进制并不会改变电路中的 bit，只是让人更容易观察边界、掩码和地址。

\begin{workedexample}
将 \(\texttt{0x2D7}\) 转换为十进制和二进制：

\[
\texttt{0x2D7}=2\times16^2+13\times16+7
=512+208+7=727.
\]

逐个十六进制数字展开为四位：

\[
\texttt{2}\rightarrow0010,\quad
\texttt{D}\rightarrow1101,\quad
\texttt{7}\rightarrow0111,
\]

所以 \(\texttt{0x2D7}=0010\,1101\,0111_2\)。前导 0 是否保留取决于
所讨论的宽度；在一个 12-bit 字段中应保留，在纯数学数值中可以省略。
\end{workedexample}

\topic{二进制运算不是十进制运算换一种字体}
位置权重决定进位规则。二进制每一位只有 0 和 1，因此

\[
0+0=0,\quad 0+1=1,\quad 1+1=10_2,\quad 1+1+1=11_2.
\]

最后两个式子分别表示“本位为 0/1，并向高位进 1”。硬件加法器正是逐位产生
和与进位；第三章会把这张算术表变成逻辑门。

\begin{workedexample}
计算 \(0110_2+0011_2\)：

\[
\begin{array}{r}
  0110\\
 +0011\\
 \hline
  1001
\end{array}
\]

从右向左依次为 \(0+1=1\)、\(1+1=0\) 进 1、
\(1+0+1=0\) 进 1、\(0+0+1=1\)，结果为 \(9\)。
\end{workedexample}

\topic{固定位宽会产生回绕、进位和溢出}
数学整数没有固定上限，寄存器和字段却有位宽。一个 \(n\)-bit 无符号加法只
保留结果对 \(2^n\) 的余数：

\[
\operatorname{result}=(a+b)\bmod 2^n.
\]

超出最高位的部分可由 carry 条件记录。它不是“CPU 算错”，而是接口只承诺
保存 \(n\) 位。

\begin{center}
\begin{tikzpicture}[font=\footnotesize\sffamily]
  \foreach \x in {0,...,15} {
    \coordinate (p\x) at ({90-22.5*\x}:2.2);
    \fill[BookBlue] (p\x) circle (0.045);
    \node[font=\scriptsize] at ({90-22.5*\x}:2.58) {\x};
  }
  \draw[->,BookBlue,very thick] (90:2.2) arc[start angle=90,end angle=-247.5,radius=2.2];
  \node[align=center] at (0,0) {4-bit 无符号环\\加 1 沿圆周前进};
  \node[BookRed,align=center] at (4.4,0.8) {\(15+1\rightarrow0\)\\产生进位};
  \draw[os arrow] (15:2.65) -- (3.6,0.8);
\end{tikzpicture}
\end{center}
\diagramnote{环形图只表示模 \(16\) 的编码行为，不表示实际寄存器在空间中排成
圆。}

\topic{补码让同一加法器处理正数和负数}
\(n\)-bit 二进制补码把最高位权重解释为 \(-2^{n-1}\)，其余位保持正权重：

\[
x=-b_{n-1}2^{n-1}+\sum_{i=0}^{n-2}b_i2^i.
\]

因此 8-bit 补码范围为

\[
-128\le x\le127.
\]

求 \(-x\) 的常用位操作是“逐位取反再加 1”。例如

\[
5=\texttt{0000\,0101},\qquad
-5=\texttt{1111\,1011}.
\]

两者相加得到 \(\texttt{1\,0000\,0000}\)，丢弃第九位后为 0。同一组 bit
是无符号数还是有符号数，取决于当前操作和类型，不取决于存储单元自身。

\begin{osgridlongtable}{L{0.18\linewidth}L{0.23\linewidth}L{0.26\linewidth}L{0.23\linewidth}}
4-bit 模式 & 无符号解释 & 补码解释 & 说明 \\ \hline
\endfirsthead
4-bit 模式 & 无符号解释 & 补码解释 & 说明 \\ \hline
\endhead
\texttt{0000} & 0 & 0 & 两种解释相同 \\ \hline
\texttt{0111} & 7 & 7 & 补码最大正数 \\ \hline
\texttt{1000} & 8 & \(-8\) & 补码最小负数 \\ \hline
\texttt{1111} & 15 & \(-1\) & 所有位为 1 \\ \hline
\end{osgridlongtable}

\topic{carry 与 signed overflow 回答不同问题}
carry 表示无符号结果超出位宽；signed overflow 表示补码真值超出有符号范围。
4-bit 中 \(15+1\) 产生 carry，但按补码解释是 \(-1+1=0\)，没有 signed
overflow；\(7+1\) 没有最高位 carry，却从 \(0111\) 变成 \(1000\)，即
\(7\) 变成 \(-8\)，产生 signed overflow。

\begin{workedexample}
对 8-bit 运算 \(\texttt{0x7F}+\texttt{0x01}\)：

\[
\texttt{0111\,1111}+\texttt{0000\,0001}
=\texttt{1000\,0000}.
\]

无符号解释是 \(127+1=128\)，合法且无 carry；补码解释期望
\(127+1=128\)，超出最大值 127，因此发生 signed overflow。条件标志必须与
程序想采用的解释配套使用。
\end{workedexample}

\topic{按位运算是在位置集合上操作}
AND、OR、XOR 和 NOT 对同一位置独立运算。常量
\term{掩码}{mask}用 1 表示“关心或修改该位”，用 0 表示“保持或忽略”。

\begin{osgridlongtable}{L{0.19\linewidth}L{0.25\linewidth}L{0.46\linewidth}}
目的 & 表达式 & 位级效果 \\ \hline
\endfirsthead
目的 & 表达式 & 位级效果 \\ \hline
\endhead
读取字段 & \(x\land mask\) & 掩码为 0 的位置被清零 \\ \hline
设置若干位 & \(x\lor mask\) & 掩码为 1 的位置变为 1 \\ \hline
清除若干位 & \(x\land\lnot mask\) & 掩码为 1 的位置变为 0 \\ \hline
翻转若干位 & \(x\oplus mask\) & 掩码为 1 的位置取反 \\ \hline
测试第 \(k\) 位 & \(x\land(1\ll k)\) & 结果非零表示该位为 1 \\ \hline
\end{osgridlongtable}

\begin{workedexample}
设备状态寄存器为 \(\texttt{1011\,0100}_2\)，bit 2 表示 ready，bit 5 表示
error。掩码分别为

\[
1\ll2=\texttt{0000\,0100},\qquad
1\ll5=\texttt{0010\,0000}.
\]

\[
\texttt{1011\,0100}\land\texttt{0000\,0100}
=\texttt{0000\,0100}\ne0,
\]

\[
\texttt{1011\,0100}\land\texttt{0010\,0000}
=\texttt{0010\,0000}\ne0.
\]

两个状态都已置位。这里的 bit 编号从最低有效位开始为 0；若手册采用不同编号
约定，必须先明确。
\end{workedexample}

\topic{移位既可能表示乘除，也可能只是重排位置}
对不溢出的无符号数，左移 \(k\) 位相当于乘 \(2^k\)，逻辑右移相当于向下取整
除 \(2^k\)。但移位也常用于构造字段和拆地址，此时不应把它强行理解成算术。

\[
\texttt{0011\,0101}\ll2
=\texttt{1101\,0100}\quad\text{（只保留 8 bit）}.
\]

有符号右移是否复制符号位属于语言和指令语义；不能只画箭头猜测。移位计数、
位宽和被移走的位都是接口的一部分。

\topic{字节序回答多字节值怎样排进连续地址}
假设 32-bit 值为 \(\texttt{0x12345678}\)，占用地址 \(A\) 到 \(A+3\)。
x86 使用 little-endian，最低有效 byte 放在最低地址：

\begin{center}
\begin{tikzpicture}[font=\footnotesize\sffamily]
  \foreach \i/\addr/\byte in {
    0/\(A\)/\texttt{78},
    1/\(A+1\)/\texttt{56},
    2/\(A+2\)/\texttt{34},
    3/\(A+3\)/\texttt{12}} {
    \node[draw=BookBlue,fill=BookCodeBg,minimum width=26mm,minimum height=9mm]
      (b\i) at ({3.0*\i},0) {\byte};
    \node[below=2mm of b\i] {\addr};
  }
  \draw[decorate,decoration={brace,amplitude=5pt}]
    (-1.3,0.65) -- (10.3,0.65)
    node[midway,above=6pt]{同一个 32-bit 值：\texttt{0x12345678}};
  \node[BookTeal,align=center] at (1.4,-1.25) {低地址\\低有效 byte};
  \node[BookAmber,align=center] at (7.6,-1.25) {高地址\\高有效 byte};
\end{tikzpicture}
\end{center}

字节序不改变一个 byte 内的 bit 编号，也不改变寄存器中这个数的数学值。只有
当多字节对象被拆成有地址的 byte、发送到协议或写入磁盘格式时，排列规则才
显现。

\topic{按地址逐字节重建值}
若内存从低到高包含
\(\texttt{EF BE AD DE}\)，按 x86 little-endian 读取 32-bit 无符号值：

\[
\texttt{0xDEADBEEF}.
\]

推导顺序是

\[
\texttt{EF}\times2^0+
\texttt{BE}\times2^8+
\texttt{AD}\times2^{16}+
\texttt{DE}\times2^{24}.
\]

磁盘格式、网络协议和设备寄存器可能指定不同字节序，所以可靠代码要显式编码
和解码，不能把宿主机内存布局原样复制后假定到处一致。

\topic{对齐、填充与半开区间是布局推理的最低工具}
\term{对齐}{alignment}要求对象起始地址是某个粒度的整数倍。例如
8-byte 对齐意味着

\[
address\bmod8=0.
\]

结构体字段之间可能插入 padding，以满足后续字段对齐；结构末尾也可能填充，
使数组中下一个元素仍然对齐。因此“字段大小之和”等于结构大小并不总成立。

\begin{center}
\begin{tikzpicture}[font=\scriptsize\sffamily]
  \foreach \i in {0,...,15} {
    \draw[BookRule] (\i*0.72,0) rectangle ++(0.72,0.72);
    \node[below] at (\i*0.72+0.36,0) {\i};
  }
  \fill[BookBlue!22] (0,0) rectangle (0.72,0.72);
  \node at (0.36,0.36) {A};
  \fill[BookRed!12] (0.72,0) rectangle (2.88,0.72);
  \node at (1.8,0.36) {padding};
  \fill[BookTeal!22] (2.88,0) rectangle (5.76,0.72);
  \node at (4.32,0.36) {32-bit B};
  \fill[BookAmber!22] (5.76,0) rectangle (11.52,0.72);
  \node at (8.64,0.36) {64-bit C};
  \node[align=center] at (5.76,-1.0) {示意布局：具体 ABI 可能不同，必须由
  \texttt{sizeof}/\texttt{alignof} 与二进制契约确认};
\end{tikzpicture}
\end{center}

\term{半开区间}{half-open interval}写作 \([begin,end)\)：包含 begin，不包含
end。长度直接为

\[
length=end-begin.
\]

相邻区间 \([a,b)\) 与 \([b,c)\) 不重叠也不留缝。地址、缓冲区和磁盘扇区范围
都适合用这种表示。

\begin{workedexample}
缓冲区从地址 \(\texttt{0x1000}\) 开始，长度
\(\texttt{0x180}\) byte，则区间为

\[
[\texttt{0x1000},\texttt{0x1180}).
\]

最后一个有效 byte 地址是 \(\texttt{0x117F}\)，不是
\(\texttt{0x1180}\)。检查访问 \([p,p+n)\) 时还必须防止
\(p+n\) 发生固定位宽溢出。
\end{workedexample}

\begin{failurebox}
“这个数放得进 64 bit”不代表它是合法地址；“结构字段看起来连续”不代表没有
填充；“读四个 byte”不代表可以在任意未对齐地址安全读取；“内存里看到
\texttt{78 56 34 12}”也不代表人类书写值是 \texttt{0x78563412}。每个结论都
需要位宽、类型、字节序、对齐和地址空间上下文。
\end{failurebox}


\section{状态编码、可靠传输与性能量}

一个开关可以处于断开或闭合；一根满足电气条件的数字线可以被接收器判为低或
高；一个触发器可以保存两个稳定状态。只要两个状态能够可靠区分，就可以约定
它们分别表示 0 和 1。

\begin{pdfdiagram}
  \node[os hardware] (physical0) {物理状态 A\\低电压范围};
  \node[os hardware,below=of physical0] (physical1) {物理状态 B\\高电压范围};
  \node[os state,right=22mm of physical0,yshift=-10mm] (bit) {逻辑解释\\0 或 1};
  \node[os block,right=22mm of bit] (meaning) {协议赋义\\数据、地址或控制};
  \draw[os arrow] (physical0) -- (bit);
  \draw[os arrow] (physical1) -- (bit);
  \draw[os arrow] (bit) -- (meaning);
\end{pdfdiagram}
\diagramnote{同一物理高电平可以在某根线上表示数据 1，在低有效复位线上表示
“没有复位”。意义来自协议和上下文，不来自电压本身。}

\topic{一个 bit 与多个 bit}
一个 bit 有 \(2\) 种状态；\(n\) 个相互独立的 bit 有

\[
2^n
\]

种组合。8 bit 有 \(2^8=256\) 种组合，通常组成一个 byte（字节）。一个
无符号 \(n\)-bit 数能表示

\[
0\le x\le 2^n-1.
\]

所以 8-bit 无符号数范围是 0--255，16-bit 是 0--65535。范围上限要减 1，
因为状态总数包含 0。

\begin{workedexample}
某设备寄存器有 5 bit 用于队列编号。最多可以编码
\(2^5=32\) 个编号，即 0--31。若需要单独保留一个编码表示“无效”，可用于
真实队列的编码只剩 31 个。硬件字段宽度给出的是编码容量，不自动等于可用
对象数量。
\end{workedexample}

\topic{编码不是对象本身}
字符 `A' 可以采用 ASCII 编码为 \(\texttt{0x41}\)；像素颜色可以编码为三组
通道值；机器指令也编码为字节序列。相同字节只有在解释规则相同时才有相同
含义。

\begin{osgridlongtable}{L{0.18\linewidth}L{0.25\linewidth}L{0.47\linewidth}}
位模式 & 解释规则 & 可能含义 \\ \hline
\endfirsthead
位模式 & 解释规则 & 可能含义 \\ \hline
\endhead
\texttt{0x41} & ASCII & 字符 `A' \\ \hline
\texttt{0x41} & 无符号整数 & 十进制 65 \\ \hline
\texttt{0x41} & x86 指令流的一部分 & 可能是 REX 前缀，含义取决于后续字节与模式 \\ \hline
\texttt{0x41} & 设备寄存器字段 & 由该设备数据手册定义 \\ \hline
\end{osgridlongtable}

\topic{数据量、地址数量与传输速度必须分开}
bit（b）和 byte（B）大小写不同：

\[
1\ \mathrm{B}=8\ \mathrm{b}.
\]

SI 前缀按 1000 进位，IEC 二进制前缀按 1024 进位：

\[
1\ \mathrm{MB}=10^6\ \mathrm{B},\qquad
1\ \mathrm{MiB}=2^{20}\ \mathrm{B}=1\,048\,576\ \mathrm{B}.
\]

“64-bit CPU”通常描述通用寄存器和指令集能力，不等于机器一定安装
\(2^{64}\) byte RAM，也不等于一次总线事务总会传 64 bit。寄存器宽度、
虚拟地址有效位、物理地址有效位、数据总线宽度和内存容量是不同规格。

\begin{workedexample}
一条标称 \(\SI{1}{Gb/s}\) 的链路，每秒原始 bit 数为 \(10^9\)。即使暂时
忽略编码、协议头、间隙和重传，换算为 byte 也只有

\[
\frac{10^9}{8}=125\,000\,000\ \mathrm{B/s}=125\ \mathrm{MB/s}.
\]

实际应用吞吐一定不大于原始线速换算值。看到“1 G”时必须先确认是 bit/s、
byte/s、十进制容量还是频率。
\end{workedexample}

\topic{从“能表示多少状态”到“怎样可靠传送状态”}
若 \(n\) 个 bit 可以独立取 0 或 1，就有 \(2^n\) 种模式。但模式数量只是编码
空间，协议还可能保留部分模式、增加校验，或让某些组合非法。可靠通信必须
同时回答：

\begin{enumerate}[leftmargin=2.2em]
  \item 一帧从哪里开始、在哪里结束；
  \item 每组 bit 按什么顺序解释；
  \item 发送和采样在什么时刻发生；
  \item 接收方怎样发现错误；
  \item 出错后丢弃、重试还是带错继续。
\end{enumerate}

\topic{并行与串行只描述组织方式，不直接决定快慢}
并行接口同一时刻可在多根线上传多个 bit，但线间偏斜、连接器引脚数和串扰随
宽度增长。串行接口用更少信号线按时间连续发送 bit，却可以通过更高信号速率、
差分传输、嵌入时钟和多条 lane 获得很高吞吐。

\begin{pdfdiagram}[node distance=8mm and 13mm]
  \node[os hardware,minimum width=25mm] (pa) {并行发送端};
  \node[os hardware,right=57mm of pa,minimum width=25mm] (pb) {并行接收端};
  \foreach \y/\label in {7/D0,4/D1,1/D2,-2/D3} {
    \draw[os arrow] ([yshift=\y mm]pa.east) -- node[above,font=\scriptsize]{\label}
      ([yshift=\y mm]pb.west);
  }
  \node[BookMuted] at ($(pa)!0.5!(pb)+(0,-12mm)$) {四根数据线在同一采样时刻组成一个 4-bit 字};

  \node[os hardware,below=25mm of pa,minimum width=25mm] (sa) {串行发送端\\并行转串行};
  \node[os hardware,right=57mm of sa,minimum width=25mm] (sb) {串行接收端\\串行转并行};
  \draw[os bus] (sa) -- node[above]{一根数据通道：\(D_0,D_1,D_2,D_3,\ldots\)} (sb);
  \node[BookMuted] at ($(sa)!0.5!(sb)+(0,-9mm)$) {同一通道在多个符号时刻依次承载 bit};
\end{pdfdiagram}

吞吐近似为

\[
\text{有效吞吐}
=\text{每秒符号数}
\times\text{每符号有效 bit}
\times\text{lane 数}
\times\text{协议效率}.
\]

最后一个因子小于等于 1，因为帧头、校验、空闲、编码和重传都占用传输机会。

\topic{时钟同步、异步和时钟恢复}
同步接口额外传送或共享时钟，接收方在规定边沿采样。异步 UART 不传连续时钟，
而是双方约定接近的波特率，接收方用起始位定位一帧并在预计中心采样。高速串行
接口常把时钟信息编码进数据变化，由接收端时钟数据恢复电路重建采样节拍。

\begin{center}
\begin{tikzpicture}[x=0.82cm,y=0.72cm,font=\footnotesize\sffamily]
  \draw[->,BookMuted] (0,0) -- (14.5,0) node[right]{时间};
  \node[left] at (0,3.4) {UART};
  \draw[BookBlue,very thick]
    (0,3.8) -- (1,3.8) -- (1,3.0)
    -- (2.5,3.0) -- (2.5,3.8)
    -- (4,3.8) -- (4,3.0)
    -- (5.5,3.0) -- (5.5,3.8)
    -- (7,3.8) -- (7,3.0)
    -- (8.5,3.0) -- (8.5,3.8)
    -- (10,3.8) -- (10,3.0)
    -- (11.5,3.0) -- (11.5,3.8) -- (14,3.8);
  \foreach \x/\label in {1.75/start,3.25/D0,4.75/D1,6.25/D2,7.75/D3,9.25/D4,10.75/D5,12.25/stop} {
    \draw[dashed,BookRule] (\x,2.65) -- (\x,4.15);
    \node[below,rotate=35,anchor=east,font=\scriptsize] at (\x,2.6) {\label};
    \fill[BookRed] (\x,3.4) circle (0.055);
  }
  \node[BookRed,align=center] at (7.5,1.2)
    {接收器在每个预计 bit 中心采样；\\边沿只帮助定位，不代表采样点本身};
\end{tikzpicture}
\end{center}

\topic{冗余让错误变得可观察}
若所有编码都用于有效数据，接收一个模式时无法仅凭该模式知道传输是否出错。
加入冗余可以检测部分错误。

\topic{奇偶校验}
偶校验增加一个 bit，使整组 1 的数量为偶数。例如数据
\(\texttt{1011001}\) 中有 4 个 1，偶校验位为 0；若传输中恰有一个 bit
翻转，接收方会看到奇数个 1。

奇偶校验能检测所有奇数个 bit 翻转，却不能检测偶数个翻转，也不能告诉哪一位
错了。错误检测能力必须与故障模型一起描述。

\topic{校验和、CRC 与纠错码}
\begin{osgridlongtable}{L{0.18\linewidth}L{0.32\linewidth}L{0.40\linewidth}}
机制 & 基本思路 & 能做什么与不能做什么 \\ \hline
\endfirsthead
机制 & 基本思路 & 能做什么与不能做什么 \\ \hline
\endhead
重复发送 & 多次发送同一信息并比较 & 简单但效率低；多个副本同错时仍可能失败 \\ \hline
奇偶校验 & 约束 1 的个数奇偶性 & 检测奇数个翻转，不能定位错误 \\ \hline
加法校验和 & 按固定规则累加数据块 & 便于软件实现，对某些重排或抵消错误较弱 \\ \hline
CRC & 把 bit 串视为多项式并计算余数 & 对突发错误有明确检测能力，仍不是密码学认证 \\ \hline
ECC & 加入足够冗余以定位部分错误 & 可纠正规定数量错误，不能保证任意损坏可恢复 \\ \hline
哈希/MAC & 摘要或带密钥认证 & 面向完整性甚至攻击者，但不能替代介质错误恢复 \\ \hline
\end{osgridlongtable}

\begin{fromzerobox}
“发现损坏”“纠正损坏”“证明发送者身份”和“故障后恢复旧版本”是四个不同
目标。CRC 不能证明数据来自可信者；数字签名不能自动修复被破坏的磁盘块；
冗余副本若同时被错误软件覆盖，也不再是可恢复备份。
\end{fromzerobox}

\topic{延迟、带宽、吞吐和并发不能互相替代}
\term{延迟}{latency}描述一次操作从开始到完成要多久；\term{带宽}{bandwidth}
描述单位时间可搬运多少；\term{吞吐}{throughput}描述系统实际完成多少工作；
\term{并发度}{concurrency}描述同时在途的操作数量。

若单次访问往返延迟为 \(L\)，每个请求大小为 \(S\)，同时最多有 \(N\) 个请求
在途，则仅由延迟隐藏能力给出的吞吐上界近似为

\[
\text{throughput}\le\frac{N S}{L},
\]

同时还不能超过物理链路带宽和目标设备处理能力。

\begin{workedexample}
磁盘读取每个请求 \(\SI{128}{KiB}\)，平均完成延迟
\(\SI{2}{ms}\)，最多 8 个请求并发。理想上界：

\[
\frac{8\times\SI{128}{KiB}}{\SI{2}{ms}}
=\SI{512}{MiB/s}.
\]

若接口实际只有 \(\SI{300}{MB/s}\)，最终上界取两者较小值，而且协议、排队和
软件开销还会继续降低实际吞吐。只降低单次延迟或只提高链路带宽都未必消除
系统瓶颈。
\end{workedexample}

\topic{编码分层：字符、整数、指令和文件格式各有契约}
同一串 byte 可以跨多层解释。以终端输出字符 `A' 为例：

\begin{pdfdiagram}[node distance=7mm and 10mm]
  \node[os block,minimum width=22mm] (meaning) {人的意图\\字符 `A'};
  \node[os state,right=of meaning,minimum width=22mm] (codepoint) {字符编码\\ASCII 65};
  \node[os state,right=of codepoint,minimum width=22mm] (byte) {byte 模式\\\texttt{0x41}};
  \node[os hardware,right=of byte,minimum width=25mm] (frame) {UART 帧\\起始/数据/停止};
  \node[os hardware,right=of frame,minimum width=22mm] (voltage) {引脚波形\\高低电压随时间};
  \draw[os arrow] (meaning) -- node[above,font=\scriptsize]{编码} (codepoint);
  \draw[os arrow] (codepoint) -- node[above,font=\scriptsize]{序列化} (byte);
  \draw[os arrow] (byte) -- node[above,font=\scriptsize]{成帧} (frame);
  \draw[os arrow] (frame) -- node[above,font=\scriptsize]{驱动} (voltage);
\end{pdfdiagram}

接收方沿相反方向解帧、取 byte、按字符编码解释，最后显示字形。任一层契约
不一致都可能产生乱码，但“乱码”不能直接告诉我们错误位于字体、编码、字节序、
串口参数还是电气连接。

\begin{failurebox}
看到一串 byte 时直接宣布“这是文本”“这是指令”或“这是地址”，等于跳过解释
契约。高质量分析应先记录来源、长度、边界、字节序、版本和校验，再选择解释器；
解释失败时保留原始 byte 作为证据。
\end{failurebox}


\topic{误差、范围和边界条件}
真实测量不是无限精确。标称 \(\SI{10}{k\ohm}\)、误差 \(1\%\) 的电阻，实际
允许范围是

\[
\SI{10}{k\ohm}\times(1\pm0.01)
=\SIrange{9.9}{10.1}{k\ohm}.
\]

数字系统通过给 0 和 1 留出电压范围来容纳噪声，而不是要求某条线精确等于
\(\SI{0}{V}\) 或 \(\SI{3.3}{V}\)。软件中的范围检查也是同一种思路：
先定义允许集合，再拒绝集合外输入。

\begin{pdfdiagram}[x=1cm,y=1cm]
  \draw[->,draw=BookMuted,thick] (0,0) -- (12,0) node[right]{输入电压};
  \fill[BookTeal!15] (0,0.15) rectangle (3.3,0.65);
  \fill[BookRed!10] (3.3,0.15) rectangle (7.0,0.65);
  \fill[BookBlue!12] (7.0,0.15) rectangle (11.5,0.65);
  \node at (1.65,0.4) {可靠低};
  \node at (5.15,0.4) {未保证区};
  \node at (9.25,0.4) {可靠高};
  \draw (3.3,0.05) -- (3.3,0.75) node[above]{\(V_{IL(max)}\)};
  \draw (7.0,0.05) -- (7.0,0.75) node[above]{\(V_{IH(min)}\)};
\end{pdfdiagram}
\diagramnote{这是概念图，不给出某一具体逻辑族的阈值。实际数值必须查器件
数据手册，并同时满足输出端与输入端规格。}

\section{计算机历史怎样逼出操作系统}

“计算”早于电子计算机。人类先用手指、刻痕、算盘和纸笔保存中间结果，再用
机械装置自动完成某类运算。后来继电器、电子管、晶体管和集成电路依次改善
速度、可靠性、体积、功耗与制造规模。历史的重点不是背年份，而是理解为什么
今天的机器同时保留数字逻辑、存储程序、层次存储和兼容接口。

\begin{osgridlongtable}{L{0.16\linewidth}L{0.21\linewidth}L{0.27\linewidth}L{0.26\linewidth}}
阶段 & 主要载体 & 解决的问题 & 留给今天的思想 \\ \hline
\endfirsthead
阶段 & 主要载体 & 解决的问题 & 留给今天的思想 \\ \hline
\endhead
人工计数 & 刻痕、算盘、纸笔 & 保存数量与中间步骤 & 状态、位置计数、算法步骤 \\ \hline
机械计算 & 齿轮、连杆、打孔介质 & 自动重复确定操作 & 控制序列与可编程输入 \\ \hline
机电计算 & 继电器 & 用电控制开关，减少机械传动 & 离散逻辑与布尔电路 \\ \hline
电子管计算 & 真空电子器件 & 大幅提高开关速度 & 电子数字计算与时钟控制 \\ \hline
晶体管计算 & 半导体开关 & 降低体积、功耗并提高寿命 & 现代逻辑门的器件基础 \\ \hline
集成电路 & 单片多晶体管 & 降低连接成本，提高密度 & 芯片、标准逻辑与模块化 \\ \hline
微处理器 & 单芯片 CPU & 让通用处理器成为可复制部件 & ISA、软件生态与个人计算机 \\ \hline
现代系统 & 多核 CPU/SoC、层次存储 & 缓解速度、功耗和并行瓶颈 & 缓存、并发、虚拟化与专用加速 \\ \hline
\end{osgridlongtable}

\begin{historybox}
早期机器常通过重新接线、开关或外部介质改变任务。存储程序思想把“将执行的
步骤”也编码成存储器里的数据，CPU 因而能按地址取出指令并改变后续控制流。
这并不意味着指令和普通数据在概念上相同，而是它们可以共享同一种位存储与
传输基础设施；执行权限、地址空间和解释规则仍会区分二者。
\end{historybox}

\topic{九个转折点怎样连接到今天}
\begin{osgridlongtable}{L{0.14\linewidth}L{0.22\linewidth}L{0.28\linewidth}L{0.26\linewidth}}
大致时期 & 转折 & 当时要突破的限制 & 今天仍能看到的结果 \\ \hline
\endfirsthead
大致时期 & 转折 & 当时要突破的限制 & 今天仍能看到的结果 \\ \hline
\endhead
古代到近代 & 位置计数和机械辅助 & 人脑难以可靠保存大量中间结果 & 进位、数位权重、算法步骤 \\ \hline
19 世纪 & 打孔介质与机械可编程设想 & 每换任务都重新制造或接线 & 程序与执行机器可以分离 \\ \hline
19 世纪中叶 & 布尔代数 & 逻辑推理缺少可演算形式 & AND/OR/NOT 和逻辑化简 \\ \hline
20 世纪上半叶 & 继电器/电子开关逻辑 & 机械部件慢且磨损 & 离散开关电路和自动控制 \\ \hline
1940 年代 & 通用电子数字计算 & 专用计算流程难以复用 & 指令、寄存器和控制单元 \\ \hline
1940--1950 年代 & 存储程序体系 & 程序配置慢，控制与数据设备割裂 & CPU 按地址取指，程序可生成程序 \\ \hline
1940 年代末以后 & 晶体管替代电子管 & 体积、功耗、热和寿命限制规模 & 半导体逻辑与低电压数字系统 \\ \hline
1950--1960 年代 & 集成电路 & 分立元件和焊点数量限制可靠性 & 芯片、标准单元和大规模集成 \\ \hline
1970 年代以后 & 微处理器与个人计算机 & 通用 CPU 昂贵且难以复制部署 & ISA 生态、主板、外设和兼容历史 \\ \hline
\end{osgridlongtable}

这些转折没有一条单独的“发明日”能完整概括。例如存储程序思想来自理论、
工程项目、存储器可用性与多位研究者的共同推进；实际机器也经历“设计思想提出”
“局部实现”“稳定投入使用”不同时间。读历史时应避免把复杂演进压成某个人
某一天突然完成整台现代计算机的神话。

\topic{操作系统为什么会出现}
最早的机器资源昂贵，操作人员手工装入任务，CPU 在准备输入输出时大量空闲。
批处理系统让任务排队自动衔接；中断和 DMA 等机制让 I/O 与计算重叠；
多道程序在一个任务等待时运行另一个任务；分时系统又用短时间片给多个交互
用户提供“机器正在响应我”的体验。

\begin{osgridlongtable}{L{0.18\linewidth}L{0.29\linewidth}L{0.43\linewidth}}
阶段 & 主要问题 & 操作系统增加的责任 \\ \hline
\endfirsthead
阶段 & 主要问题 & 操作系统增加的责任 \\ \hline
\endhead
人工操作 & 装带、拨开关和记录结果占用机器 & 尚无稳定 OS，操作流程由人承担 \\ \hline
批处理 & 任务切换成本高 & 作业队列、自动装入、错误后继续下一任务 \\ \hline
多道程序 & CPU 等待慢速 I/O & 中断、内存分区、调度和设备管理 \\ \hline
分时 & 多个用户需要交互响应 & 抢占、隔离、终端、账号和共享文件 \\ \hline
虚拟内存 & 程序超过物理内存且彼此需隔离 & 地址翻译、分页、按需装入和保护 \\ \hline
Unix 式系统 & 工具需要组合，硬件接口需要统一 & 进程、fd、层次文件系统、pipe 与 Shell \\ \hline
现代系统 & 多核、网络、安全和设备规模巨大 & 并发、权限、协议栈、驱动模型和资源控制 \\ \hline
\end{osgridlongtable}

因此操作系统不是为了“让屏幕好看”才出现，而是在昂贵共享硬件上持续回答：
谁获得 CPU、谁拥有哪段内存、慢设备完成时通知谁、数据怎样命名并在故障后
存活、不可信程序之间如何隔离。本项目随后会亲手建立这些责任。

\topic{从人工操作到操作系统：问题怎样一层层逼出抽象}
历史不是一串互不相干的发明。每一层新抽象通常来自上一层机器在规模、速度或
共享方面遇到的瓶颈。

\topic{人工计算时期：算法和执行者还没有分开}
纸笔、算盘和表格已经包含计算机思想的雏形：输入、规则、中间状态和结果。
但执行规则主要保存在人的训练和文档里，速度、重复性和错误率受操作者限制。
机械计算器把某些步骤固化进齿轮，使“规则”部分转移到机器结构。

\topic{可编程介质：改变任务不必重造整台机器}
打孔卡、纸带和插线板让控制序列能够脱离部分机械结构。这里出现了一个关键
分离：

\[
\text{通用执行机构}+\text{可更换控制描述}=\text{可编程机器}.
\]

早期“程序”不一定存放在今天意义上的 RAM 中，但已经回答了为什么同一台机器
可以重复承担不同任务。

\topic{继电器和电子管：逻辑速度提高，可靠性成为新瓶颈}
继电器用电磁作用控制机械触点，容易理解和隔离，但开关慢、会磨损并产生触点
抖动。电子管取消机械触点，大幅提高速度，却需要较大功耗、体积和维护成本。
机器越大，器件数量越多，单个器件较小的失效率也会累积成系统可用性问题。

\topic{晶体管和集成电路：连接本身成为制造对象}
晶体管比电子管更小、更省电且更耐用；集成电路进一步把晶体管与互连一起制造。
真正的跨越不只是“元件变小”，还包括：

\begin{itemize}[leftmargin=2.2em]
  \item 大量连接从人工焊接变成可复制的版图；
  \item 单位逻辑功能成本下降，复杂控制器成为可能；
  \item 芯片接口可以标准化，系统可按模块组合；
  \item 规模继续扩大后，功耗、布线延迟和验证又成为新限制。
\end{itemize}

\topic{存储程序：指令成为可寻址的编码}
当指令和数据都能以编码形式进入可寻址存储，CPU 可以维护“下一条指令地址”，
顺序取指，也可以通过跳转改变控制流。编译器、装载器、调试器和操作系统由此
拥有共同基础：程序不再只是插线板状态，而是可以复制、检查、重定位和保护的
字节布局。

\begin{pdfdiagram}[node distance=7mm and 6mm]
  \node[os hardware,minimum width=25mm] (rewire) {专用/重新接线\\人工改变机器};
  \node[os hardware,right=of rewire,minimum width=25mm] (media) {外部程序介质\\控制序列可更换};
  \node[os state,right=of media,minimum width=25mm] (stored) {存储程序\\指令可寻址};
  \node[os block,right=of stored,minimum width=25mm] (software) {软件工具链\\编译、装载、调试};
  \node[os block,below=9mm of software,minimum width=28mm] (os) {操作系统\\共享、保护和持久化};
  \draw[os arrow] (rewire) -- (media);
  \draw[os arrow] (media) -- (stored);
  \draw[os arrow] (stored) -- (software);
  \draw[os arrow] (software) -- (os);
\end{pdfdiagram}

\topic{为什么 CPU 变快以后反而更需要操作系统}
早期 CPU 极其昂贵，人工装入下一任务的几分钟会浪费大量计算能力。批处理监控
程序先解决“自动接续任务”；I/O 远慢于 CPU 后，多道程序解决“一个任务等待时
让另一个任务运行”；多人交互需求又推动分时、权限和文件共享。

\begin{center}
\begin{tikzpicture}[x=0.82cm,y=0.68cm,font=\footnotesize\sffamily]
  \draw[->,BookMuted] (0,0) -- (15,0) node[right]{时间};
  \node[left] at (0,4.2) {单任务};
  \fill[BookBlue!25] (0.5,3.8) rectangle (3.2,4.5);
  \node at (1.85,4.15) {CPU 计算};
  \fill[BookRed!15] (3.2,3.8) rectangle (8.5,4.5);
  \node at (5.85,4.15) {等待 I/O，CPU 空闲};
  \fill[BookBlue!25] (8.5,3.8) rectangle (11.2,4.5);
  \node at (9.85,4.15) {继续};

  \node[left] at (0,2.4) {多道};
  \fill[BookBlue!25] (0.5,2.0) rectangle (3.2,2.7);
  \node at (1.85,2.35) {任务 A};
  \fill[BookTeal!25] (3.2,2.0) rectangle (6.8,2.7);
  \node at (5.0,2.35) {A 等待时运行 B};
  \fill[BookAmber!25] (6.8,2.0) rectangle (8.5,2.7);
  \node at (7.65,2.35) {任务 C};
  \fill[BookBlue!25] (8.5,2.0) rectangle (11.2,2.7);
  \node at (9.85,2.35) {任务 A};
  \node[align=left] at (12.8,3.15)
    {需要：中断通知\\保存/恢复现场\\内存隔离\\调度策略};
\end{tikzpicture}
\end{center}

由此可以看到，操作系统职责不是凭空设计出来的：

\begin{osgridlongtable}{L{0.20\linewidth}L{0.31\linewidth}L{0.39\linewidth}}
机器矛盾 & 若没有新机制会怎样 & 被逼出的 OS 抽象 \\ \hline
\endfirsthead
机器矛盾 & 若没有新机制会怎样 & 被逼出的 OS 抽象 \\ \hline
\endhead
CPU 快、I/O 慢 & CPU 长时间空等设备 & 中断、阻塞/唤醒、异步驱动 \\ \hline
多个任务共享 CPU & 一个任务可永久霸占执行 & 定时器、抢占和调度 \\ \hline
多个程序共享内存 & 错误地址破坏别的程序或内核 & 地址空间、页表和权限 \\ \hline
设备接口各异 & 每个程序重复操纵硬件寄存器 & 驱动、统一 I/O 与系统调用 \\ \hline
数据需要跨重启 & 内存断电后全部丢失 & 块层、文件系统、命名和恢复 \\ \hline
用户彼此不可信 & 一个程序能窃取或破坏全部状态 & 特权级、身份、访问控制和审计 \\ \hline
\end{osgridlongtable}

\topic{为什么兼容性让现代 x86 看起来“不整齐”}
成功的平台会积累大量已有软件、设备和知识。新处理器若完全抛弃旧接口，性能
也许更整齐，却会失去生态。x86 的复位模式、历史设备端口、可变长指令和多种
执行模式因此同时存在。兼容性不是单纯技术保守，而是已有投资与新设计之间的
工程权衡。

本书后续会反复采用同一历史读法：

\begin{enumerate}[leftmargin=2.2em]
  \item 当时机器有什么限制；
  \item 新机制解决了哪一个主要问题；
  \item 为兼容旧软件保留了什么；
  \item 新机制又制造了哪些复杂性；
  \item 本项目在 QEMU 平台上选择实现哪一部分；
  \item 哪些实机问题被模型暂时省略。
\end{enumerate}

\begin{historybox}
“旧”不等于“无用”，“新”也不等于“没有代价”。例如中断减少 CPU 轮询浪费，
却引入异步控制流和并发；分页提供保护与虚拟内存，却引入页表、TLB 和页故障；
缓存提高平均速度，却引入一致性与可见性问题。理解历史问题，才能理解机制为何
以今天的形状存在。
\end{historybox}

\topic{本章纸笔实验：不接触硬件也能形成可检验证据}
\begin{enumerate}[leftmargin=2.2em]
  \item 给出 \(\SI{250}{MHz}\)，计算周期并写出完整单位消去过程；
  \item 把 \(\texttt{0xB37A}\) 展开成 16 bit，再按 little-endian 写出四个
        地址上的 byte；
  \item 对 8-bit 运算 \(\texttt{0xFF}+1\) 分别写出无符号与补码解释；
  \item 为一个含 ready、error 和 3-bit queue id 的寄存器设计位布局和掩码；
  \item 画出字符 `A' 从人的含义到 UART 引脚波形的五层编码链；
  \item 选择“中断”“分页”或“文件系统”，用“旧限制—新机制—新代价”写出
        一页因果说明。
\end{enumerate}

验收不只看答案，还要检查是否写明位宽、单位、字节序、地址区间和解释契约。


\section{模型、证据与从零推理}

模型主动忽略细节。地铁图忽略真实弯曲距离，却保留换乘关系；指令集模型忽略
晶体管布局，却保留软件可见寄存器和指令结果；QEMU TCG 在另一种宿主 CPU
上复现目标 x86-64 CPU 和选定设备的可观察行为。

\begin{pdfdiagram}
  \node[os hardware] (board) {实体机器\\电源、芯片、铜线};
  \node[os state,right=of board] (arch) {架构/设备契约\\寄存器、时序、结果};
  \node[os block,right=of arch] (qemu) {QEMU 模型\\在宿主机上执行};
  \node[os block,below=of arch] (book) {本书推理\\只在声明边界内成立};
  \draw[os arrow] (board) -- node[above]{呈现} (arch);
  \draw[os arrow] (qemu) -- node[above]{复现} (arch);
  \draw[os arrow] (arch) -- (book);
\end{pdfdiagram}

QEMU 能证明项目 ROM 字节是否被模拟 CPU 从复位地址取到，也能证明对 UART
端口的访问是否遵循设备模型；它不能证明真实板上的电源纹波、焊点、走线阻抗
或散热满足要求。本书后续每一张图都会标明它描述的是逻辑关系、架构事务还是
实体电气连接。

\topic{从零推理的六步纸笔法}
以后遇到任何机制，都可以在纸上写出六栏：

\begin{enumerate}[leftmargin=2.2em]
  \item \textbf{对象：}有哪些元件、寄存器、内存区或软件对象？
  \item \textbf{量与单位：}宽度、地址、时间、电压、容量分别是什么？
  \item \textbf{初始状态：}上电、复位或函数进入时已经保证什么？
  \item \textbf{转换规则：}哪条线、指令或函数让状态发生变化？
  \item \textbf{边界：}允许范围、提交点和禁止状态是什么？
  \item \textbf{证据：}用计算、测量、反汇编、日志或反例怎样推翻错误答案？
\end{enumerate}
